% SPDX-License-Identifier: CC-BY-4.0
% Copyright (c) 2026 Martin Schröder <info@swedishembedded.com>

\section{The learning loop}\label{sec:loop}\label{app:loop}

Three projects divide the work, and the division is enforced by a layering file and a gate rather than
by convention: \textbf{Sven} runs agents and keeps the record of what an agent did, \textbf{Brain}
computes (inference, training, checkpoints, adapters, serving), and \textbf{Splinter} decides what is
learned, which experience qualifies as training data, what to train and whether to release
it.\cl{S05,S01,S02} \Cref{fig:arch} shows one turn. A source is captured with its provenance; a model
proposes tasks and code admits or refuses each for a named reason; an agent attempts them in a confined
copy of the tree; verifiers that the policy cannot reach keep what they confirm and an exposure audit
withholds attempts that touched hidden material; verified experience is projected into training records
with whole held-out groups; Brain trains an adapter or every weight from the current champion; and a
release gate admits the candidate only if four checks on material it was not trained on pass
(improvement by a paired sign test, retention of earlier releases, an anchor of general tasks, and
agreement of the served model with the trained one). Releases are immutable, an alias moves, and
rollback moves it back.\cl{A03,V02,L08,G01,X02,X06} Every stage is specified, with its ledger, in
\cref{app:system}; this section is what the experiments below need.

In the experiments of \cref{sec:rsi} the loop is reduced to one task kind whose answers an interpreter
fixes, the verifier is the interpreter, and the release gate is replaced by the learner's own dev gate
and, for the registered decision, by a separate examiner program. In \cref{sec:coding-summary} the full
loop runs on repository tasks.

\begin{figure}[tbp]
\centering
\begin{tikzpicture}[fig]
\tikzset{cell/.style={text width=29mm}}
\def\dx{39mm}
% row 1, left to right
\node[data, cell] (src) at (0,0) {\textbf{source}:\\document, repository,\\command or session};
\node[proc=okSky!22, cell] (task) at (\dx,0) {\textbf{tasks}: a model\\proposes, code\\admits (16 kinds)};
\node[proc=okOrange!22, cell] (roll) at (2*\dx,0) {\textbf{attempts}: a Sven\\agent in a confined\\copy of the tree};
\node[verif, cell] (ver) at (3*\dx,0) {\textbf{verifiers} and\\exposure audit\\(13 verifier kinds)};
% row 2, right to left
\node[data, cell] (data) at (3*\dx,-31mm) {\textbf{dataset}: views of\\verified experience,\\held-out groups};
\node[proc=okPurple!22, cell] (train) at (2*\dx,-31mm) {\textbf{training} in Brain:\\LoRA, DPO or full,\\from the champion};
\node[verif, cell] (gate) at (\dx,-31mm) {\textbf{release gate}:\\improvement,\\retention, anchor,\\serve};
\node[data, cell] (rel) at (0,-31mm) {\textbf{release}:\\written once; an\\alias moves; rollback};
% flows
\draw[flow] (src) -- (task);
\draw[flow] (task) -- (roll);
\draw[flow] (roll) -- node[lab, right, pos=0.5] {} (ver);
\draw[flow] (ver) -- node[lab, right] {verified\\experience} (data);
\draw[flow] (data) -- (train);
\draw[flow] (train) -- node[lab, above] {} (gate);
\draw[flow] (gate) -- (rel);
\draw[ctrl] (rel.north) -- node[lab, right, align=left] {next turn runs on\\the released policy} ++(0,9mm) -| ($(roll.south)+(-6mm,0)$);
% what each stage refuses
\draw[refuse] (src.north) -- ++(0,5mm) node[reflabel, above, align=center] {over 1 MiB, not text,\\dirty tree, no user step};
\draw[refuse] (task.north) -- ++(0,5mm) node[reflabel, above, align=center] {21 reasons: ungrounded,\\duplicate, contradiction};
\draw[refuse] (roll.north) -- ++(0,5mm) node[reflabel, above, align=center] {outside the tree,\\network: kernel-refused};
\draw[refuse] (ver.north) -- ++(0,5mm) node[reflabel, above, align=center] {failed, exposed,\\undecided: withheld};
\draw[refuse] (data.south) -- ++(0,-5mm) node[reflabel, below, align=center] {compacted, over the\\block, near-duplicate};
\draw[refuse] (gate.south) -- ++(0,-5mm) node[reflabel, below, align=center] {any check failed\\or not measured};
\draw[refuse] (rel.south) -- ++(0,-5mm) node[reflabel, below, align=center] {exists already,\\terms forbid it};
% legend of the projects
\node[proc=okSky!22, text width=14mm, minimum height=5mm, inner sep=1.5pt] (k1) at (0.6*\dx,-59mm) {Splinter};
\node[proc=okOrange!22, text width=14mm, minimum height=5mm, inner sep=1.5pt, right=4mm of k1] (k2) {Sven};
\node[proc=okPurple!22, text width=14mm, minimum height=5mm, inner sep=1.5pt, right=4mm of k2] (k3) {Brain};
\node[verif, text width=14mm, minimum height=5mm, inner sep=1.5pt, right=4mm of k3] (k4) {verifier};
\node[data, text width=14mm, minimum height=5mm, inner sep=1.5pt, right=4mm of k4] (k5) {record};
\end{tikzpicture}
\caption{The self-learning loop, read clockwise from the source. Every stage keeps what passes and
refuses the rest with a named reason (dotted arrows; \cref{sec:intake,sec:admission,sec:verify,sec:dataset,sec:gate}
give the full lists). Splinter owns tasks, verification, datasets, the gate and releases; Sven runs
the attempts; Brain trains and serves. The dashed arrow is the next turn: attempts run on the policy
the release alias points at.}\label{fig:arch}
\end{figure}
